%% sections/results.tex

\chapter{Resultados}
\label{chap:results}

Neste capítulo, apresentamos os resultados das quatro análises conduzidas pela equipe. Cada seção aborda uma dimensão diferente da pergunta de pesquisa.

% ══════════════════════════════════════════════════════════════════════════
\section{Tendências Temporais}
\label{sec:temporal}

A análise temporal, conduzida por Gabriel, buscou entender como a confiabilidade dos lançamentos da SpaceX mudou ao longo do período de 2006 a 2022.

\subsection{Evolução da Taxa de Sucesso}

A trajetória da taxa de sucesso ao longo dos anos pode ser vista na \autoref{fig:gabriel_success}. Nos anos iniciais (2006--2008), quando a SpaceX ainda operava o Falcon 1, os resultados eram instáveis — a taxa de sucesso variou entre 0\% e 50\%. A partir de 2010, com a entrada do Falcon 9, o cenário mudou bastante: a taxa subiu para 100\% e permaneceu acima de 85\% em todos os anos seguintes. De 2017 em diante, todos os lançamentos foram bem-sucedidos.

\begin{figure}[H]
    \centering
    \includegraphics[width=0.85\textwidth]{gabriel_success_rate_per_year.png}
    \caption{Taxa de sucesso dos lançamentos da SpaceX por ano (2006--2022), com intervalos de confiança de 95\%.}
    \label{fig:gabriel_success}
\end{figure}

\subsection{Crescimento Operacional}

Outro aspecto relevante é o volume de lançamentos. A SpaceX saiu de 1--2 lançamentos por ano entre 2006 e 2010 para 43 lançamentos em 2022 (\autoref{fig:gabriel_launches}). O que chama atenção é que esse aumento de escala não veio acompanhado de perda de qualidade — a confiabilidade se manteve alta durante todo o período de crescimento.

\begin{figure}[H]
    \centering
    \includegraphics[width=0.85\textwidth]{gabriel_launches_per_year.png}
    \caption{Número de lançamentos da SpaceX por ano (2006--2022).}
    \label{fig:gabriel_launches}
\end{figure}

\subsection{Adoção da Reutilização}

A prática de reutilizar boosters foi se consolidando ao longo do tempo. Antes de 2016, nenhum booster havia sido reutilizado. Em 2020, praticamente 100\% dos boosters utilizados já tinham voos anteriores (\autoref{fig:gabriel_reuse}). Esse crescimento acompanhou o período de maior confiabilidade da empresa.

\begin{figure}[H]
    \centering
    \includegraphics[width=0.85\textwidth]{gabriel_reuse_rate_per_year.png}
    \caption{Taxa de reutilização de boosters por ano.}
    \label{fig:gabriel_reuse}
\end{figure}

% ══════════════════════════════════════════════════════════════════════════
\section{Impacto da Reutilização}
\label{sec:reusability}

Kaio foi direto na pergunta central do projeto: será que reutilizar um booster compromete o resultado da missão? Para responder, ele separou os registros em dois grupos — virgens e reutilizados — e comparou os resultados.

\subsection{Comparação Global}

Os números estão na \autoref{tab:kaio_stats}. Dos 149 voos com boosters reutilizados, nenhum falhou. Já entre os 43 voos com boosters virgens, 5 resultaram em falha, o que dá uma taxa de 88,37\%. A diferença entre os dois grupos é de 11,63 pontos percentuais.

\begin{table}[H]
\centering
\caption{Taxas de sucesso comparadas entre boosters virgens e reutilizados.}
\label{tab:kaio_stats}
\begin{tabular}{lrrr}
\toprule
\textbf{Tipo} & \textbf{Quantidade} & \textbf{Sucessos} & \textbf{Taxa de Sucesso} \\
\midrule
Virgem       & 43  & 38  & 88,37\% \\
Reutilizado  & 149 & 149 & 100,00\% \\
\midrule
\textbf{Total} & \textbf{192} & \textbf{187} & \textbf{97,40\%} \\
\bottomrule
\end{tabular}
\end{table}

\begin{figure}[H]
    \centering
    \includegraphics[width=0.85\textwidth]{kaio_overall_success_rate.png}
    \caption{Comparação visual das taxas de sucesso: virgens vs.\ reutilizados.}
    \label{fig:kaio_overall}
\end{figure}

\subsection{Controle Temporal}

Para verificar se essa diferença não era apenas um reflexo da época, o Kaio estratificou os dados por ano (\autoref{fig:kaio_temporal}). O que se observou é que as falhas dos boosters virgens se concentram nos primeiros anos do programa (2006--2015). A partir de 2016, quando os reutilizados começam a aparecer, ambos os grupos apresentam taxas de sucesso altas.

\begin{figure}[H]
    \centering
    \includegraphics[width=0.85\textwidth]{kaio_temporal_distribution.png}
    \caption{Taxas de sucesso ao longo do tempo, separadas por tipo de booster.}
    \label{fig:kaio_temporal}
\end{figure}

% ══════════════════════════════════════════════════════════════════════════
\section{Experiência do Booster}
\label{sec:experience}

A análise de Izac aprofundou a investigação: em vez de apenas dividir entre ``virgem'' e ``reutilizado'', ele criou faixas de experiência para avaliar se existe algum padrão de degradação conforme o booster acumula mais voos.

\subsection{Estratificação por Faixas de Experiência}

Os resultados estão na \autoref{tab:izac_stats}. O padrão é bem claro: boosters em seu primeiro voo (novatos) têm taxa de sucesso de 88,37\%, mas a partir do segundo voo a taxa sobe para 100\% e se mantém em todas as faixas seguintes — independentemente de o booster ter 2 ou 13 voos anteriores.

\begin{table}[H]
\centering
\caption{Taxa de sucesso por faixa de experiência do booster.}
\label{tab:izac_stats}
\begin{tabular}{lrr}
\toprule
\textbf{Nível de Experiência} & \textbf{N} & \textbf{Taxa de Sucesso} \\
\midrule
Novatos (0 voos anteriores)       & 43  & 88,37\% \\
Intermediários (1--3 voos)        & 50  & 100,00\% \\
Experientes (4--6 voos)           & 30  & 100,00\% \\
Veteranos (7+ voos)               & 69  & 100,00\% \\
\bottomrule
\end{tabular}
\end{table}

\begin{figure}[H]
    \centering
    \includegraphics[width=0.85\textwidth]{izac_experience_vs_success.png}
    \caption{Relação entre a experiência do booster e a taxa de sucesso.}
    \label{fig:izac_experience}
\end{figure}

\subsection{Ausência de Degradação}

Esse resultado é particularmente relevante: se houvesse desgaste acumulado nos boosters, seria de esperar que a taxa de sucesso caísse aos poucos com o aumento do número de voos. Mas não é isso que aparece nos dados. A relação se comporta mais como uma função degrau — o primeiro voo concentra todo o risco, e após esse ``batismo'', o booster mantém desempenho consistente.

% ══════════════════════════════════════════════════════════════════════════
\section{Famílias de Foguetes}
\label{sec:families}

Por fim, Jonatas comparou a confiabilidade das três famílias de foguetes operadas pela SpaceX.

\subsection{Comparação entre Famílias}

A \autoref{tab:jonatas_stats} mostra as taxas de sucesso de cada família. A diferença entre gerações é bastante visível: o Falcon 1 ficou em 40\%, o Falcon 9 em 98,9\% e o Falcon Heavy em 100\%.

\begin{table}[H]
\centering
\caption{Confiabilidade por família de foguete.}
\label{tab:jonatas_stats}
\begin{tabular}{lrrr}
\toprule
\textbf{Família} & \textbf{Total} & \textbf{Sucessos} & \textbf{Taxa de Sucesso} \\
\midrule
Falcon 1     & 5   & 2   & 40,0\% \\
Falcon 9     & 178 & 176 & 98,9\% \\
Falcon Heavy & 9   & 9   & 100,0\% \\
\bottomrule
\end{tabular}
\end{table}

\begin{figure}[H]
    \centering
    \includegraphics[width=0.85\textwidth]{jonatas_taxa_sucesso_familia.png}
    \caption{Taxa de sucesso por família de foguete.}
    \label{fig:jonatas_family}
\end{figure}

Do ponto de vista estatístico, o Falcon 9 é de longe o mais relevante pelo volume de dados disponível (178 registros). O Falcon Heavy tem resultado perfeito, mas com apenas 9 registros é difícil tirar conclusões fortes. Já o Falcon 1, operado entre 2006 e 2009, pertence a uma fase muito diferente da empresa — tanto em tecnologia quanto em experiência operacional.
